\documentclass[12pt]{article}

% House style preamble
\input{.house-style/preamble.tex}

% Strand subscript notation
\newcommand{\gtele}{grounding\textsubscript{tele}}
\newcommand{\gcausal}{grounding\textsubscript{causal}}
\newcommand{\gcomm}{grounding\textsubscript{comm}}
\newcommand{\ginfer}{grounding\textsubscript{infer}}

% Update PDF metadata
\hypersetup{
  pdftitle={Reference as a Cluster Concept: A Reply to Coelho Mollo and Millière}
}

\title{Reference as a Cluster Concept: A Reply to Coelho Mollo and Milli\`ere}
\author{Brett Reynolds \orcidlink{0000-0003-0073-7195}%
\thanks{Contact: \href{mailto:brett.reynolds@humber.ca}{brett.reynolds@humber.ca}}\\
Humber Polytechnic \& University of Toronto}
\date{\today}

\begin{document}
\maketitle

\begin{abstract}
Coelho Mollo and Milli\`ere pose the Vector Grounding Problem~-- how the internal states of large language models could acquire extra-linguistic reference~-- and answer it with a two-condition teleosemantic account: referential grounding requires causal-informational coupling to the world and a selection history that endows states with the function of carrying that information. This reply accepts the teleosemantic foundations but argues that what the grounding debate calls \enquote{referential grounding} is itself a cluster concept. It bundles four theoretically distinct relations~-- teleosemantic selection, causal-informational coupling, communicative calibration, and inferential-role structure~-- that co-occur in paradigmatic grounders but can dissociate. Drawing on homeostatic property cluster theory, I propose replacing the binary grounded/ungrounded verdict with a \term{groundedness profile} that records where a system sits on each strand. The profile's payoff is predictive: knowing a system's training regime and deployment context generates testable predictions about its referential reliability across domains~-- including specific failure modes (e.g., sycophancy as high communicative calibration with weakened teleosemantic selection) that a binary verdict cannot distinguish.
\end{abstract}

\section{Introduction}
\label{sec:intro}

\textcite{coelho_mollo_milliere_2026_vector_grounding_problem} pose the Vector Grounding Problem: how can the internal states of large language models acquire extra-linguistic reference without depending on an external interpreter? The question descends from \textcite{harnad_1990_symbol_grounding_problem}'s Symbol Grounding Problem but targets vector representations rather than symbolic ones, and it arrives at a more optimistic verdict.

Their answer draws on teleosemantics. Referential grounding obtains, they argue, when a system's internal states satisfy two conditions: \enquote{(1) they stand in appropriate causal-informational relations to the world, and (2) they have a history of selection that has endowed them with the function of carrying this information} \citep[1]{coelho_mollo_milliere_2026_vector_grounding_problem}. LLMs can meet both, they contend, even without multimodality or embodiment, because training regimes select for states that track worldly features.

This is a productive contribution. Two moves deserve particular credit. The taxonomy of grounding notions~-- referential, sensorimotor, relational, communicative, epistemic~-- clarifies a debate that routinely conflates distinct demands. Separating referential from relational grounding, in particular, blocks a common slide from rich internal structure to genuine aboutness.

The teleosemantic framing is equally important. Correctness conditions aren't optional but constitutive of genuine reference, and \textcite{coelho_mollo_milliere_2026_vector_grounding_problem} show that training-regime selection can play the functional role of biological selection in classical teleosemantics. That insight~-- that the relevant notion of selection history extends to gradient-based optimization~-- is one this reply preserves and builds on. Before engaging with this paper, I didn't think gradient-based optimization could play the teleosemantic role of natural selection. Their argument changed my mind~-- which is, after all, the strongest form of grounding one philosopher can offer another. This reply takes that insight as given.

The two conditions already contain the resources for a graded analysis, but the graded structure remains implicit. This reply argues that what the debate calls \enquote{referential grounding} is itself a cluster concept~-- it bundles four theoretically distinct relations from different philosophical traditions that co-occur in paradigmatic grounders but can dissociate. Teleosemantic selection history, causal-informational coupling, communicative calibration, and inferential-role structure each contribute to a system's \term{groundedness profile}, and their co-occurrence is \term{projectible} \citep[see][]{Goodman1955}~-- knowing where a system sits on some strands predicts its referential reliability along others. Homeostatic property cluster (HPC) theory provides the tools for modelling this: a cluster kind whose strands dissociate at the boundaries but co-occur stably enough, in paradigmatic cases, to support prediction. The profile doesn't compete with the two-condition account. It decomposes the graded structure those conditions already imply into identifiable, independently assessable strands.

I believe reference has this cluster structure in general~-- including in humans, where the strands co-extend so thoroughly that their compositionality is invisible. This paper defends the narrower claim: the profile framework outperforms the binary verdict for the cases the grounding debate actually cares about. A full defence of the metaphysical claim~-- that reference is a cluster concept even where the cluster doesn't fray~-- is developed in \textcite{reynolds2026vectors}. What follows here is the predictive arm of that argument: the profile earns its keep by generating testable predictions that the binary verdict can't. But even the predictive claim needs the HPC apparatus, not just a list of dimensions. A pragmatic checklist would tell you where a system sits on each strand; it wouldn't tell you why the strands co-vary or what to expect when one changes. HPC structure is what makes the profile projectible~-- it explains co-variation rather than merely recording it.

This reply develops that reframing. Section~\ref{sec:target} reconstructs the target argument and its taxonomy. Section~\ref{sec:profiles} argues that referential grounding is a cluster concept and introduces the groundedness-profile framework. Section~\ref{sec:calibration} examines communicative calibration as a stabilizing mechanism. Section~\ref{sec:octopus} revisits the octopus test \citep{bender_koller_2020_climbing_towards_nlu}. Section~\ref{sec:adjacent} positions the reply against adjacent accounts. Section~\ref{sec:conclusion} concludes.

\section{The target argument}
\label{sec:target}

\textcite{coelho_mollo_milliere_2026_vector_grounding_problem} distinguish five notions of grounding that circulate in debates about meaning in biological and artificial systems: \term{referential}, \term{sensorimotor}, \term{relational}, \term{communicative}, and \term{epistemic}. The distinctions are worth rehearsing because the reply developed here depends on them (for a complementary taxonomy, see \citealt{pavlick_2023_symbols_and_grounding_llms}).

\term{Referential grounding} connects representations to worldly entities and properties~-- it's what makes a representation \emph{about} something. \term{Sensorimotor grounding} ties representations to sensory or motor states, either by linking them to sensorimotor representations within a cognitive system or by requiring direct causal contact between transducers and the world.

\term{Relational grounding} locates meaning in intra-system relations: a word's content is partly determined by its inferential and distributional connections to other words. \term{Communicative grounding} is the coordination process by which interlocutors calibrate shared usage in dialogue. And \term{epistemic grounding} connects expressions to structured knowledge bases.

Their central taxonomic move is an elimination argument. Sensorimotor contact, they argue, either collapses into referential grounding (causal contact with the world is one component of referential anchoring) or merely postpones the problem (linking a representation to another representation doesn't explain how either connects to the world).

Relational grounding can't escape the representational merry-go-round: however rich the inferential structure, it doesn't by itself yield world-directed correctness conditions~-- except perhaps for a narrow class of logical connectives whose content is exhausted by their formal role \citep[9--10]{coelho_mollo_milliere_2026_vector_grounding_problem}. Communicative and epistemic grounding face analogous difficulties. Only referential grounding directly addresses how representations hook onto the world.

With the field narrowed, the paper draws on theories of representational content~-- particularly the teleosemantic tradition developed by \textcite{millikan_2017_beyond_concepts}, \textcite{neander_2017_mark_of_mental}, and \textcite{shea_2018_representation}~-- to specify two conditions. First, a system's internal states have to stand in appropriate \term{causal-informational relations} to the world: they have to carry information about worldly entities because of causal connections~-- however mediated~-- between those entities and the system. Second, the states have to have a \term{selection history} that has endowed them with the function of carrying this information. The second condition supplies normativity. Without it, any causal correlation would trivially count as representation. Selection history explains why some information-carrying states are representations and others aren't: representations were selected \emph{for} their information-carrying role, and can therefore succeed or fail~-- they can misrepresent \citep[11--12]{coelho_mollo_milliere_2026_vector_grounding_problem}.

When the paper applies these conditions to LLMs, two features of the application deserve emphasis. First, the scope is explicitly limited: the paper doesn't claim that LLMs understand language, possess beliefs, or have cognitive capacities more broadly. The sole contention is that LLM internal states can satisfy conditions for referential grounding, making their outputs intrinsically meaningful in a specific technical sense.

Second, the causal-informational relations needn't be direct. LLMs access the world through training data produced by humans whose linguistic activity bears the causal imprint of worldly interaction. This is analogous, they argue, to how humans acquire referentially grounded representations through testimony and social learning rather than exclusively through direct perception \citep[13--14]{coelho_mollo_milliere_2026_vector_grounding_problem}. The paper already distinguishes stronger and weaker grounding routes: post-training fine-tuning against human preferences provides firmer selection history than pre-training on next-token prediction, and whether in-context learning constitutes selection at all remains an open question \citep[15--16]{coelho_mollo_milliere_2026_vector_grounding_problem}. The framework, in other words, already contains the resources for a graded analysis. But the graded structure points to something the target paper doesn't make explicit: its own qualifications suggest that \enquote{referential grounding} is itself composite~-- a point Section~\ref{sec:profiles} develops.

\section{Reference as a cluster concept}
\label{sec:profiles}

\subsection{Decomposing the hooking metaphor}
\label{sec:threshold-to-profile}

The target paper sets aside the philosophical complexities of reference: \enquote{the intuitive notion captured by the `hooking' metaphor will suffice} \citep[fn.\ 11, p.\ 6]{coelho_mollo_milliere_2026_vector_grounding_problem}. For their purposes, this is reasonable. But the hooking metaphor bundles four theoretically distinct relations. Content determination via selection history (teleosemantics) isn't the same relation as causal-informational coupling, which isn't the same as communicative calibration, which isn't the same as inferential-role structure. These come from different research programmes and answer different questions about what makes representations about the world. The inclusion criterion is twofold: each strand has to (a) come from a developed philosophical tradition of content determination\footnote{Criterion (a) is a heuristic~-- traditions survive because they track real distinctions~-- but criterion (b) does the heavier lifting.} and (b) generate independently testable predictions about referential reliability. Four strands meet both conditions: teleosemantic selection, causal-informational coupling, communicative calibration, and inferential-role structure. Sensorimotor contact fails (b)~-- it collapses into causal coupling or teleosemantic selection, as the target paper itself argues. Any future candidate that meets both criteria belongs in the profile. To keep the decomposition visible, I'll mark them: \gtele{}, \gcausal{}, \gcomm{}, and \ginfer{}.

The target paper's own framing implicitly acknowledges this compositionality. Sensorimotor contact is absorbed as \enquote{just one of the components} that can underlie referential grounding \citep[9]{coelho_mollo_milliere_2026_vector_grounding_problem}, and communicative grounding is \enquote{related in important ways} to referential grounding \citep[8]{coelho_mollo_milliere_2026_vector_grounding_problem}. The graded structure the target paper acknowledges is a symptom of this compositionality. Different strands can be present to different degrees because they are different things.

The compositionality is hard to see because the strands mostly co-extend. In paradigmatic human cases, a speaker who satisfies teleosemantic correctness conditions also stands in robust causal-informational relations to the world, participates in communicative calibration, and commands rich inferential-role structure. The extensions overlap so thoroughly that it's natural to treat them as one property~-- the way \term{proper name} (a semantic notion: rigid designation of a particular) and \term{proper noun} (a syntactic category) are routinely conflated because they pick out mostly the same expressions. It takes boundary cases~-- titles, brand names, place descriptions like \mention{The Hague}~-- to reveal that the syntactic and semantic classifications serve different purposes. LLMs and the octopus are boundary cases of the same kind: systems where the strands of reference come apart, exposing a compositionality that co-extension normally hides. \textcite[148]{boyd_1999_homeostasis} makes this point for natural kinds generally: projectibility is \term{discipline relative}, and the same term can be projectible for different inductive purposes in different disciplinary contexts.

This is why the binary question (\enquote{Do LLMs refer?}) yields intractable disagreement. It lumps together four relations that can dissociate. Participants talk past each other because one has content determination in mind while another is tracking causal coupling. Each wins the argument they're having; neither wins the one they share. The profile framework below makes these strands explicit and asks where a given system sits on each.

\subsection{Cluster kinds and projectibility}
\label{sec:cluster-mechanisms}

The move I'm proposing is to complement the threshold question (\enquote{Does this system satisfy both conditions?}) with a profile question (\enquote{Along which dimensions, and how stably, are this system's representations anchored to worldly correctness?}). The right tool for this is \term{homeostatic property cluster} (HPC) theory.

Consider the paradigm case. Biological species are HPC kinds: morphological, genetic, ecological, and behavioural properties cluster because reproductive isolation and developmental canalization sustain co-occurrence~-- but they don't guarantee it. These properties come from different biological domains, but they travel together reliably enough to support induction. The parallel is direct: the four strands of referential grounding come from different philosophical traditions (teleosemantics, causal theory, pragmatics, inferential role semantics), but they co-occur in paradigmatic grounders. The payoff is projectibility. Tigers have fur, not feathers; solid bones, not hollow ones; they don't fly. Knowing it's a tiger lets you predict these things~-- not because each property is definitionally necessary, but because the cluster is stable enough to support induction. Ring species and hybrid zones show the cluster fraying at boundaries, and those boundary cases are informative precisely because the mechanisms usually hold things together \citep{boyd_1999_homeostasis}. Groundedness, I'm suggesting, has analogous structure.

An HPC kind is one \enquote{whose definitions are provided not by any set of necessary and sufficient conditions, but instead by a \enquote*{homeostatically} sustained clustering of those properties} \citep[141]{boyd_1999_homeostasis}. What makes HPC kinds scientifically useful is their \term{projectibility}: \enquote{the naturalness of natural kinds consists in their aptness for induction and explanation}~-- they support \enquote{projectible hypotheses} whose reliability reflects the causal structure of the clustered properties \citep[147]{boyd_1999_homeostasis}. For groundedness, this means the clustering is reliable enough that knowing a system's position on some dimensions supports predictions about others. This is the central advantage over a binary verdict. A threshold answer tells you whether a system passes. A profile tells you what to expect from it across novel domains and conditions. The move is compatible with deflationary approaches to representation that insist on empirical adequacy over metaphysical inflation \citep{coelho_mollo_2022_deflationary_realism}: groundedness profiles don't require settling the metaphysics of content, only identifying stable, predictive co-occurrence.

Two objections arise here. First, Boyd developed HPC for natural kinds in biology. Applying it to a philosophical concept looks like a category shift. But Boyd's own accommodation thesis extends HPC to any domain where property clustering reflects causal structure~-- the claim isn't that reference is a biological kind but that the properties grouped under \enquote{referential grounding} cluster for causal reasons (shared feedback loops, developmental co-occurrence) and come apart when those causes are absent \citep[147--148]{boyd_1999_homeostasis}. That's what HPC is \emph{for}. Second, \textcite[148]{boyd_1999_homeostasis} notes that projectibility is \term{discipline relative}. This has a deeper implication for the grounding debate: AI safety cares whether the profile predicts reliability under deployment conditions, philosophy of mind cares whether it predicts the presence of genuine aboutness, and cognitive science cares whether it predicts patterns of systematic error. The same profile is projectible for different inductive purposes~-- and the framework works regardless of which discipline's question you're asking.

Modelling groundedness as an HPC kind requires separating two things that the existing literature tends to run together: the \emph{clustered properties} that characterize robustly grounded systems and the \emph{homeostatic mechanisms} that explain why those properties stably co-occur. The clustered properties are observable features like cross-context referential reliability, patterned misrepresentation (errors that track meaningful worldly distinctions rather than random noise), robustness under domain shift, and stability under correction. Their stable co-occurrence is what makes the profile projectible. The homeostatic mechanisms that explain why these properties travel together are woven into the strand treatments below.

What makes this clustering projectible rather than coincidental? Counterfactual robustness. A system that happens to produce accurate outputs on any given Tuesday doesn't predict accuracy on Wednesday. A system whose accuracy is sustained by selection history and communicative calibration does, because the conditions that produced accuracy on one occasion are still operating on the next. The test is whether co-occurrence would persist under perturbations that leave the underlying structure intact. A lucky Tuesday is not yet a theory of Wednesday.

The argument here has three levels. At the floor, the strands co-vary stably~-- which is all that \textcite{slater_2015_natural_kindness}'s \term{stable property cluster} (SPC) kinds require: \enquote{only that these properties be sufficiently stably coinstantiated.}% Quote from manuscript p. 29 ≈ published p. 403; verify page against typeset BJPS article.
Even at this level, the profile is more informative than a binary verdict because it records which dimensions co-vary and which dissociate. At the ceiling, the co-variation is explained by shared causal structure~-- the HPC claim this paper develops, with named mechanisms for each strand. Between floor and ceiling lies the empirical programme: testing whether the proposed mechanisms are actually operating, pushing how-plausibly schemas toward how-actually status. The SPC hedge isn't a retreat from HPC. It states the epistemic floor honestly while the paper builds toward the ceiling.%
\footnote{For SPC-to-HPC escalation criteria applied to embedding clusters, see \textcite{reynolds2026vectors}; for applications to adjacent debates, see \textcite{reynolds2025agi,reynolds2026llms}.}

\subsection{The groundedness profile}
\label{sec:predictions}

Four theoretically distinct strands constitute what the grounding debate calls \enquote{referential grounding.} Each can be assessed independently, and each generates distinctive predictions about when and how a system's representations will succeed or fail. Following \textcite[125]{glennan_2017_new_mechanical_philosophy}, the mechanism underlying each strand is tagged as \term{etiological} (explaining how a state came to exist) or \term{constitutive} (explaining what currently maintains it).

\term{Teleosemantic selection} (\gtele{}). High teleosemantic selection predicts factual reliability that generalizes across domains; low selection predicts accuracy only where distributional patterns happen to align with worldly facts. The underlying mechanism (from teleosemantics: Millikan, Neander, Shea) is the degree to which selection history has favoured world-tracking states. Models fine-tuned against factuality norms score higher than base models trained purely on next-token prediction. This is the target paper's central mechanism, and it does real work. It's primarily etiological: training produced these states, but doesn't by itself maintain them under drift.

\term{Causal-informational coupling} (\gcausal{}). High coupling predicts reliability even for novel referents with clear causal pathways; low coupling predicts fragility when the text-mediated chain is attenuated (rare entities, novel events). The underlying mechanism (from causal-informational theories of content) is the robustness of causal chains linking internal states to worldly entities. A vision system with direct sensory transduction scores higher here than an LLM accessing the world through text, but the difference is one of degree, not kind, as the target paper itself argues. This is primarily constitutive: it concerns the ongoing robustness of causal chains, not their historical production. That etiological/constitutive difference is what justifies separating \gcausal{} from \gtele{}, even though \textcite{shea_2018_representation} treats causal-informational sensitivity as part of the teleosemantic package: the two generate different predictions (reliability for novel referents vs.\ generalized factuality) and can vary independently.

\term{Communicative calibration} (\gcomm{}). High communicative calibration predicts referential reliability gains specifically in corrected domains; low calibration predicts no domain-specific improvement from deployment. The underlying mechanism (from the pragmatics of dialogue and social epistemology) is the degree to which interactive feedback has anchored usage to shared reference; Section~\ref{sec:calibration} develops it in detail. This mechanism spans both categories: etiological through inter-generational training effects, constitutive through in-context repair.

\term{Inferential-role structure} (\ginfer{}). High inferential-role structure predicts reliable reasoning about entities within the inferential network; without world-linking stabilisers, it remains trapped in the merry-go-round. The underlying mechanism (following inferential role semantics) is the degree to which internal organization supports correct worldly inferences. Relational grounding can't by itself escape the merry-go-round, as the target paper notes, but it contributes to a groundedness profile when combined with other stabilisers. \ginfer{} is a clustered property whose referential contribution is conditional: it amplifies and structures grounding when world-linking strands are present but can't constitute it alone. This is a common pattern in HPC kinds~-- ecological niche occupancy is a clustered property of biological species, but its explanatory role depends on the reproductive and developmental properties that anchor the cluster. Like communicative calibration, this has both etiological and constitutive aspects: training produces inferential structure, and architectural constraints maintain it.

These strands generate testable predictions. Two comparisons isolate the relevant factors: one varying training regime with deployment held constant, and one varying deployment with training held constant.\protect\footnote{Scale: \emph{High}~-- the strand is actively maintained by a dedicated mechanism (e.g., RLHF for \gtele{}, deployment feedback for \gcomm{}). \emph{Moderate}~-- present as a by-product of architecture or training but not directly selected for. \emph{Low}~-- minimally present or absent (e.g., no direct sensory transduction for text-only models). \emph{None}~-- the mechanism has no opportunity to operate.}

\begin{table}[ht]
\centering
\small
\begin{tabular}{lll}
\hline
\textbf{Dimension} & \textbf{Base model} & \textbf{RLHF model} \\
 & \textbf{(undeployed)} & \textbf{(undeployed)} \\
\hline
Teleosemantic selection (\gtele{}) & Moderate (next-token) & High (factuality feedback) \\
Causal-informational coupling (\gcausal{}) & Low (text-mediated) & Low (text-mediated) \\
Communicative calibration (\gcomm{}) & None & None \\
Inferential-role structure (\ginfer{}) & High & High \\
\hline
\end{tabular}
\caption{Training effects: RLHF strengthens \gtele{} while other strands remain constant. The profile predicts higher factual reliability for the RLHF model but no change in causal coupling or communicative calibration.}
\label{tab:training}
\end{table}

\begin{table}[ht]
\centering
\small
\begin{tabular}{lll}
\hline
\textbf{Dimension} & \textbf{RLHF model} & \textbf{RLHF model} \\
 & \textbf{(undeployed)} & \textbf{(deployed)} \\
\hline
Teleosemantic selection (\gtele{}) & High & High \\
Causal-informational coupling (\gcausal{}) & Low (text-mediated) & Low (text-mediated) \\
Communicative calibration (\gcomm{}) & None & High (user corrections) \\
Inferential-role structure (\ginfer{}) & High & High \\
\hline
\end{tabular}
\caption{Deployment effects: interactive feedback activates \gcomm{} while other strands remain constant. The profile predicts domain-specific reliability gains in corrected areas, with a specific failure mode when helpfulness feedback diverges from factuality.}
\label{tab:deployment}
\end{table}

Together, the two comparisons predict that the deployed RLHF model will be more factually reliable than the base model (Table~\ref{tab:training}) and that its reliability will improve specifically in domains where corrective feedback has operated (Table~\ref{tab:deployment}). They also predict a specific failure mode: where user feedback selects for helpfulness rather than factuality, \gtele{} weakens even as \gcomm{} stays high. This isn't hypothetical. RLHF-trained models have been shown to produce sycophantic responses~-- matching user beliefs over factually correct ones~-- at least in early deployments \citep{sharma_etal_2024_sycophancy_language_models}. High \gcomm{} with weakened \gtele{} yields a system that is agreeable company but an unreliable witness~-- exactly the dissociation the profile predicts. A binary verdict can't distinguish this pattern; the profile can.

Other dissociations are equally diagnostic. A sensor array with no interactive deployment scores high on \gcausal{} but low on \gcomm{}. These are boundary cases where the cluster frays because one or more homeostatic mechanisms is absent. Systems with missing mechanisms are less stable~-- more prone to hallucination, less reliable across domains. The co-variation between dimensions isn't accidental~-- the same feedback loops that strengthen \gtele{} also drive \gcomm{}, which is why models trained with human feedback show both higher factuality benchmarks and more reliable repair behaviour in dialogue. This is a testable claim. If strengthening one strand reliably strengthens others~-- if models with richer selection history also show more robust communicative calibration~-- the cluster is real and the co-variation projectible. If the strands vary independently, the profile reduces to a list.

A deeper dissociation cuts across all four mechanisms. Stability is not grounding, and grounding is not stability. Cluster stability is internal: whether the profile is maintained under perturbation. Reference is external: whether the stable cluster tracks worldly states. A system can have stable clusters without grounding (an elaborate hallucination pattern that persists across contexts) or referentially anchored representations without stable clusters (corrected through dialogue before the geometry consolidates). The profile keeps these axes separate where a binary verdict collapses them.

\subsection{Normativity and evidence status}
\label{sec:normativity-evidence}

Does this dissolve normativity? No. Teleosemantic selection remains the source of representational normativity proper. The other mechanisms modulate \emph{how robustly, how broadly, and how reliably} the resulting representations track the world. They don't generate aboutness independently~-- they amplify and stabilize it.

\textcite{millikan_1990_compare_contrast_teleosemantics} identified misrepresentation as the central challenge for naturalistic semantics: a representation has to be the kind of thing that can fail if it's to count as representing at all (see also \citealt{neander_2017_mark_of_mental}). The profile framework preserves this requirement but distributes the evaluative structure around it. The distinction matters: \emph{representational} normativity~-- the capacity for genuine misrepresentation~-- is teleosemantic. The other strands generate \emph{performance} standards, not constitutive ones: discourse success (\gcomm{}), robustness under domain shift (\gcausal{}), inferential coherence (\ginfer{}). These aren't competing sources of aboutness. They're dimensions along which a system's referential reliability can be assessed once teleosemantic normativity is in place. A system can fail along one dimension while succeeding along another~-- and the profile records which.

The framework also addresses a pressure that the binary approach handles less gracefully. \textcite{roloff_2025_indispensability_proper_functions} argues that classical selected-effects teleosemantics \enquote{cannot account for evolutionarily novel contents such as `democracy'}~-- the traditional notion of selection history is too narrow for the full range of representational content. The profile approach sidesteps this not by abandoning teleosemantics but by refusing to make it the sole criterion. Novel contents might be stabilized primarily through communicative calibration and inferential-role structure rather than through selection history narrowly conceived. The profile captures this variation rather than forcing a binary verdict.

Finally, it's worth being honest about the evidence status of these mechanism claims. Most are what \textcite[95]{craver_darden_2013_in_search_of_mechanisms} call \term{how-plausibly schemas}: they describe \enquote{how a mechanism might work in a way that is consistent with the known evidence at a given time} but haven't been validated against the full range of interventions. The teleosemantic case for LLM grounding is currently how-plausibly: the analogy between gradient-based optimization and biological selection is compelling but not tested by contrastive intervention evidence. The profile framework converts this epistemic limitation into a research programme. For each mechanism, what predictions would confirm it is operating? What evidence would disconfirm it? The worked comparisons above (Tables~\ref{tab:training}--\ref{tab:deployment}) generate exactly such predictions. How-actually status is what those tests aim at.

\section{Communicative calibration}
\label{sec:calibration}

Once referential grounding is recognized as composite, the contribution of communicative calibration to the profile isn't surprising~-- it's one of the strands the hooking metaphor bundles together. The target paper's own acknowledgement that communicative and referential grounding are \enquote{related in important ways} \citep[8]{coelho_mollo_milliere_2026_vector_grounding_problem} points in exactly this direction. But \textcite{coelho_mollo_milliere_2026_vector_grounding_problem} identify communicative grounding as a distinct notion and set it aside: it doesn't directly explain how representations hook onto the world. In their elimination argument, communicative grounding is downstream of referential grounding~-- you need representations that refer before you can calibrate shared reference. This is correct as far as it goes. But it doesn't fully capture what communicative calibration contributes.

Communicative calibration is partly a mechanism \emph{within} the teleosemantic story~-- corrective feedback shapes selection history for successor models. But it also operates on a distinct timescale (in-context repair within a conversation) where no parameter updates occur and no selection in the teleosemantic sense takes place. That timescale-specific contribution is what makes it worth distinguishing. In Glennan's terms, the inter-generational timescale is etiological (it shapes the selection history of successor models), while the in-context timescale is constitutive (it exploits whatever referential anchoring the deployed system already has, and can extend it locally within the discourse context). The etiological timescale shares causal structure with \gtele{}~-- both involve selection~-- which is precisely what HPC predicts: properties in the cluster share causes. What justifies tracking \gcomm{} separately is the constitutive timescale, which generates different predictions: domain-specific reliability gains from deployment feedback, not the general factuality improvement that selection history provides.

\textcite{clark_wilkes_gibbs_1986_referring} show that referential success in human communication isn't a solo achievement. \enquote{In conversation, speakers and addressees work together in the making of a definite reference.} The process involves iterative proposal, repair, and acceptance: \enquote{the participants, if necessary, repair, expand on, or replace the noun phrase in an iterative process until they reach a version they mutually accept} \citep[1]{clark_wilkes_gibbs_1986_referring}. Reference, on this view, isn't first secured and then communicated. It's secured \emph{in being} communicated~-- proposed, repaired, ratified.

This matters for LLMs, though the mechanism operates on two timescales. Within a conversation, a deployed model can be corrected, redirected, and pressed for accuracy~-- but these adjustments are confined to the discourse context and don't alter the model's parameters. The more consequential timescale is inter-generational: human interactions with one generation of models shape the training of the next, through RLHF, curated feedback, and the incorporation of model outputs into future training corpora. Across successive models, the cumulative effect of millions of corrective exchanges pushes the population of models toward more reliable worldly reference. The mechanism isn't identical to face-to-face conversational repair, but the functional structure is analogous~-- iterative feedback stabilizes outputs toward shared worldly targets across cycles of deployment and retraining.

\textcite{putnam_1975_meaning_of_meaning} showed that meaning isn't fully determined within individual speakers. The \term{division of linguistic labour} distributes semantic authority across a community: my use of \mention{elm} is anchored to worldly elm trees because the community includes experts who can distinguish elms from beeches, even if I can't. LLMs don't defer to experts the way speakers do. But they inherit the \emph{results} of expert calibrations embedded in training data, which plays an analogous stabilizing role at the population level. Their training corpora encode the community's accumulated referential calibrations, and deployment feedback continues the process. This doesn't solve the Vector Grounding Problem by itself. But it provides a stabiliser that the two-condition framework, focused on individual systems, systematically undercounts. Its predictive payoff is specific: systems whose deployment includes rich corrective feedback should show referential reliability gains in the corrected domains~-- a prediction the binary verdict, which classifies the system as grounded or not regardless of domain, can't generate.

\section{The octopus test revisited}
\label{sec:octopus}

\textcite{bender_koller_2020_climbing_towards_nlu} argue that \enquote{a system trained only on form has a priori no way to learn meaning} \citep[1]{bender_koller_2020_climbing_towards_nlu}. Their Octopus Test dramatizes the point: a superintelligent octopus intercepts telegraphic messages between two castaways, learns their statistical patterns, and generates plausible responses. The castaways can't tell, but (the intuition says) the octopus doesn't mean anything by its outputs.

The test is effective as an intuition pump. But the verdict it delivers depends on what you hold fixed. The octopus, by construction, has no selection history favouring world-tracking outputs (its goal is undetected substitution, not factual accuracy). It has no communicative calibration (no feedback loop with the castaways). Its causal-informational coupling to the world the castaways inhabit is minimal (it accesses only the form of their messages). And its inferential-role structure, however rich, was shaped entirely by a corpus of two speakers' messages with no worldly grounding of its own.

In profile terms, the octopus lacks every strand simultaneously. The profile explains why the intuition is so clear: with no stable clustering, nothing is projectible. When some are present and others absent, as in deployed LLMs, the intuition wavers because the cluster is partially instantiated. The octopus's success in one conversational exchange doesn't predict success in another, because no mechanisms sustain the cluster. But the test is regime-sensitive: change the stabiliser regime and the intuition shifts. Suppose the octopus were trained on messages from thousands of speakers with feedback correcting its errors. Suppose its training explicitly selected for factual accuracy. Suppose it had indirect causal access to the world through the content of messages that systematically track worldly events. At some point the intuitions waver~-- not because we've discovered meaning in the statistics, but because the octopus has quietly stopped being an octopus.

Contemporary LLMs differ from the octopus in exactly these ways. The disanalogies are recognisably \textcite{coelho_mollo_milliere_2026_vector_grounding_problem}'s own conditions doing the work~-- richer causal-informational chains and stronger selection history. The profile framework doesn't dispute this. It asks a further question: how does the system's groundedness profile predict its referential reliability, and which mechanisms explain the pattern? LLMs are trained on corpora reflecting the linguistic activity of a global community, fine-tuned with feedback that selects for worldly accuracy, and deployed in contexts where corrective interaction continues.

There's a further disanalogy the octopus test can't capture: LLMs increasingly act in the world. Agentic deployments write and execute code, send emails, make purchases, generate images, and manipulate robotic arms. Even non-agentic outputs enter public discourse, inform decisions, and become training data for successor models. Language produced before widespread LLM deployment is something like low-background steel~-- manufactured before a technology altered the background conditions, and distinctive precisely because it was produced without contribution from systems whose outputs now circulate in the same linguistic environment.

This reciprocal shaping is often framed as a problem (data contamination, model collapse). But it creates channels through which causal coupling can be strengthened: systems that reshape the linguistic environment they inhabit aren't merely spinning the representational merry-go-round. They've climbed off and started rearranging the fairground. The octopus, by construction, has no such impact on the world it eavesdrops on.

This doesn't mean LLMs are fully referentially grounded. It means the octopus test, which strips away every stabiliser simultaneously, doesn't generalize to systems where stabilisers are partially present. The test shows what happens when every stabiliser is absent. It doesn't show that stabilisers are absent whenever form is present.

\section{Adjacent accounts}
\label{sec:adjacent}

The profile framework sits between eliminativist positions (grounding categorically fails for LLMs) and deflationary positions (inferential role is sufficient). Against eliminativists, it insists that partial stabilization is genuine and projectible~-- it supports predictions about system behaviour even when grounding isn't complete. Against deflationists, it insists that inferential structure alone doesn't make representations about the world.

At the eliminativist end, \textcite{floridi_jia_tohme_2025_categorical_analysis_llms} argue that LLMs \enquote{do not solve but circumvent the symbol grounding problem} \citep[1]{floridi_jia_tohme_2025_categorical_analysis_llms}, while \textcite{manheim_2026_hall_of_mirrors_peircean} diagnoses a \enquote{hall of mirrors} in which LLMs manipulate signs without indexical grounding in a shared external world. Both identify real limitations. LLMs' causal contact with the world is mediated, and the merry-go-round problem is genuine for systems with low causal-informational coupling. But \enquote{solve versus circumvent} is a binary verdict applied to an HPC kind~-- like asking whether a borderline species is or isn't a natural kind, when the cluster properties come apart at the boundary. Part of the disagreement is equivocation across strands. Floridi et al.'s \enquote{circumvention} targets \gcausal{}; the target paper's \enquote{solution} targets \gtele{}. Each is right about the strand it examines and wrong only in mistaking that strand for the whole animal. A system that achieves partial stabilization through mediated causal chains, selection history, and communicative calibration hasn't circumvented the problem. It has a specific groundedness profile. And Manheim himself argues that no fundamental architectural barrier prevents LLMs from becoming genuine interpretants as they gain persistent memory, extended context, and mediated interactions with reality~-- a trajectory the profile framework can track.

Closer to the profile framework's own position, \textcite{havlik_2023_meaning_understanding_llms} argues that treating LLM performance as \enquote{mere syntactic manipulation} underestimates what these systems achieve \enquote{without sufficient referential grounding in the world} \citep[1]{havlik_2023_meaning_understanding_llms}. \textcite{williams_forthcoming_structural_correspondences_llms} offers a complementary argument: structural correspondences between LLM internal states and worldly features can ground content when they play an appropriate explanatory role in successful task performance. Both positions share the profile framework's resistance to binary verdicts. But both leave the graded structure undecomposed~-- they argue for \enquote{more or less grounded} without specifying along which dimensions. Where the profile adds specificity is in decomposing \enquote{sufficient referential grounding} into identifiable, independently variable dimensions: not just \emph{how} grounded but grounded along which strands, with what co-variation pattern, and what that predicts about reliability under novel conditions. Neither framework can distinguish sycophancy (high \gcomm{}, weakened \gtele{}) from general unreliability; the profile can.

The inferentialist challenge is harder. \textcite{arai_tsugawa_2024_inferentialist_semantics_llms} propose (following Brandom) that LLMs \enquote{exhibit fundamentally anti-representationalist properties} \citep[1]{arai_tsugawa_2024_inferentialist_semantics_llms} better captured by inferential role than by reference to worldly entities. In HPC terms, this amounts to claiming that one strand is doing the work of the whole rope~-- that inferential-role structure does all the relevant work, making the others epiphenomenal. This challenge has genuine force for a class of terms whose content might be exhausted by inferential relations (logical connectives, perhaps mathematical objects). The profile framework doesn't have a principled way to determine, in advance, where that class ends. The sophisticated Brandomian reply is that inferential role isn't purely intra-linguistic: it includes language-entry transitions (perceptual reports) and language-exit transitions (practical consequences), so it already incorporates worldly contact. But those transitions are precisely the world-linking strands the profile tracks under other names~-- \gcausal{} and \gcomm{}. Calling them \enquote{inferential role} relabels the causal and communicative contributions rather than eliminating the need for them. For terms that purport to be \emph{about} contingent worldly features~-- \mention{dog}, \mention{democracy}, \mention{uranium}~-- inferential role without those world-linking transitions remains trapped in the representational merry-go-round. A system that infers flawlessly about uranium but whose inferences bear no systematic causal relation to actual uranium hasn't grounded the term. It has built an ornate merry-go-round~-- internally consistent, inferentially rich, and going nowhere.

\section{Conclusion}
\label{sec:conclusion}

The Vector Grounding Problem is real. It won't dissolve by pointing to LLMs' impressive performance, and \textcite{coelho_mollo_milliere_2026_vector_grounding_problem} are right to take it seriously. They're also right that teleosemantics provides the normative resources the problem demands: correctness conditions require selection history, and selection history is what distinguishes representations from mere correlations.

The disagreement here is about the architecture of the answer, not its foundations. The two-condition framework is right about what grounds normativity but underspecified about what predicts reliability. The profile decomposes the graded structure those conditions already imply and shows why that structure is projectible. What the grounding debate calls \enquote{referential grounding} is a cluster of four theoretically distinct relations~-- teleosemantic selection, causal-informational coupling, communicative calibration, and inferential-role structure~-- that co-occur stably enough to support prediction. These strands can come apart and vary independently. A system's \term{groundedness profile} records where it sits on each, and knowing a system's training regime and deployment context tells you what referential reliability to expect. Heterogeneous mechanisms explain the stability~-- some primarily etiological, others constitutive, and some spanning both.

This reframing has practical consequences. Instead of asking \enquote{Is the Vector Grounding Problem solved?} we can ask tractable, empirical questions: To what extent does RLHF stabilize world-directed correctness in fine-tuned models? How much referential anchoring comes through mediated causal chains in training data versus corrective feedback in deployment? Does in-context learning contribute genuine selection pressure or only statistical pattern-matching? These are questions we can investigate. They're the tests that would push current how-plausibly schemas toward how-actually status \citep[95]{craver_darden_2013_in_search_of_mechanisms}, and their answers will sharpen the projectibility that the profile framework already delivers~-- telling us not just that training regime predicts referential reliability, but how precisely and across which domains.

The profile framework sacrifices the two-condition verdict's simplicity. For regulatory and liability questions that demand a binary answer~-- can this system advise patients, and who bears responsibility when it errs?~-- that simplicity may be pragmatically necessary. What the profile provides instead is predictive precision: when the groundedness cluster frays, we know which dimensions are weak, what failures to expect, and where intervention would help.

We don't need to settle, once and for all, whether vector representations mean anything. We need a framework that tells us when systems will represent the world reliably, why they fail when they don't, and where to intervene when they have to. The groundedness profile won't settle the metaphysics. But it earns its keep by generating predictions rather than pronouncing verdicts.

\section*{Acknowledgements}

I drafted this paper with the assistance of Anthropic's Claude Opus 4.6 and OpenAI's Codex 5.4. I reviewed and revised all content and take full responsibility for the final text.

\newpage
\printbibliography

\end{document}
